\chapter*{序}
\addcontentsline{toc}{chapter}{序}

本书是为那些读过一次微积分而又想多学一点数学（特别是现代数学）的读者写的。写这样一本书的出发点是以下三种考虑：

首先，微积分的“基本问题”或者说“基本矛盾”是什么？长期以来，许多人认为是 $\varepsilon$--$\delta$。对此我有些怀疑。初学微积分在 $\varepsilon$--$\delta$ 问题上感到困难。但是即令掌握了它，也不能说就懂得了微积分。再说，所谓“基本问题”或“基本矛盾”，按习惯的说法是指贯穿始终，带动乃至决定微积分的主要内容的问题或矛盾。一门学问是否一定需要有一个或几个主要矛盾，借以展开这门学问，恐怕是可疑的事。可是我想，科学的目的在于探讨宇宙的规律，而从古希腊以来，就认为这种规律乃是数学的规律。于是随着人类社会的发展，科学面临的重大问题不断改变，科学以及作为它的不可少部分的数学就从一个阶段发展到另一个阶段。例如在资本主义出现的时代，了解天体运动的规律恐怕是人类面临的重大问题。说应用，它带动了动力学、航海学……；说思想，它促成了一次人类的思想大解放，使哥白尼、伽利略成了思想解放的巨人。正是在这个背景下，出现了牛顿、莱布尼茨……以及微积分。不妨说，这构成了一个研究宇宙规律，掌握乃至应用这些规律的一个大平台。牛顿等人的理论，有自己的问题和对象，有自己的方法，给了我们一个研究和处理问题的框架。其中一个核心问题是如何处理无穷小量。因此要有 $\varepsilon$--$\delta$。随着社会的发展，例如电磁现象、物质的分子构造又相继摆到我们面前。科学和数学又有了新问题和新方法。例如处理电磁理论，单只是“无穷小量分析”（这里借用了自欧拉时代就沿用的名称）就不够用了。例如我们不得不讨论向量、张量、外微分形式；不得不考虑坐标系（参考系）的变化；不得不研究许多本质上属于拓扑学的问题。我想把这一“套”东西当作一个新框架来看待。这样来写一本书，我们就必然涉及许多物理问题。数学既不是它的“加工订货”的产物，也不简单是工具或语言，而是与物理学以不同角度，用不同方法，但是又互相携手，共同研究大自然。我想，应该使许多大学生知道这是多么吸引人的事业，并且希望他们中间有些人能走上这条道路，这是第一点。

既然是形成了一个大平台，人类就在这个大平台上建筑了许多大厦，它们是紧密联系的。例如自从有了微积分的同时也就有了微分方程（常的和偏的）；在研究实变量函数的同时，也在研究复变量的函数。数学发展的历史不是先发展大学一二年级的课程，再发展三四年级的课程。我们不应该过分苛责目前大学里的课程安排。例如常微分方程与偏微分方程，问题不同，方法不同，可说是自成体系，如果按历史发展先后来安排，你可能会使学生一头雾水，不知所以。这样做十分不利于科学发展，而现在的安排基本上是合理的。但是也产生了一个问题，即各门课程互相孤立。从20世纪80年代起根据中法两国政府的协议，在武汉大学数学系办了一个试验班。由法国政府派数学教师在武大按法国的办法教学，以期我们能汲取对我国教育事业有益的经验。这个班给人印象甚深的是：一二年级只有一门数学课——高等数学，其内容包括了我国大学数学系一二年级全部数学基础课。此外则是物理学、计算机。到了高年级，有几位法国同行建议只要几门课就够了，而且举一些例，有微分学与复分析（是指一些常微分方程、微分流形、外微分形式等等）；测度论、概率论与泛函分析；李群与微分几何；代数方面可以讲一点表示论（群论基础知识已经在一二年级讲了）。此外可以有一些选修课。这些建议自然有一些随意性，而且有一些是日常“闲谈”讲的话，算不得教学计划。但是他们中大多数人是这个看法。问他们为什么，回答是：“数学就是统一的，应该让学生有一个统一的认识”。而且他们对我国“分工过细”常表示很不习惯。我在写这本书时，也试着在这方面作一些努力。但是这样就要回答一个问题：例如本书涉及一些所谓实分析知识，是否可以用它来取代实分析课程？我想是不行的。因为既然已分划出这样一个分支，总有自己的道理：有自己的问题，自己的方法，想要掌握它，必须系统地下功夫。所以我不幻想本书能取代其它的书，而只是告诉读者，曲径可以通幽。读者在学了微积分以后，就很容易进入某个领域，并且试着引导他去看一下这个领域的概貌。至于读者是否要进入某个领域，那是读者的事：要看有没有时间，有没有需要，而最重要的是有没有兴趣。我所能做的事无非是趁读者游兴正浓之时，告诉他还有哪些地方值得漫游，值得探胜乃至探险。一是提高他的游兴，二是如果他真的能去了，不至于过于生疏。根据同样的思想，本书可以只读其一章，可以从任一章开始，这与一般教科书也不太相同。这是第二点。

第三点是我们近年来常感到一个问题：可否让学生尽快地进入科学的前沿？另外，在研究生的教学或讨论班中又时常发现有不少内容大学本科就可以懂或者应该懂，而研究生有些弄不清的问题很可能是忘记了他在学微积分时实际上已经接触过这个问题，或者是因为变了一个样子就不认识了。科学发展太快，而学生学习年限不可能再延长。而且将来至少有一部分学生要在没有学校和老师的条件下自己去钻研新的科学知识。所以在大学中就应该为他们准备条件。结论是应该早一点让学生接触新成果。问题是能否做到。上面讲的试验班的经验给了我们一些启发：一方面是要认真地提取出新成果的实质并与传统的教学内容结合起来。否则，老是一本书一本书念下去绝不是办法。其次是老师们要“想开”一点：是不是认真教好书是教师的职责；是不是认真读好书是学生的职责。二者虽有密切关系却不是一件事。教学的某一部分内容时常并非最基本的，学生是否愿意在这里花工夫应该可以自己选择。即令听了不甚了然，至少没有什么坏处。我的一位老师当年就告诉过我，年轻时代思想最敏锐，即令一时不懂，将来再接触到会有似曾相识之感，对自己大有好处。这段话是千真万确的。这本书的内容有相当一部分是我曾以种种形式对学生们讲过的，效果不一定比讲传统的教材更差。走什么样的路，甚至数学系的学生将来是不是一定会以数学为生，这确实是学生们自己的事。“想开”了这一点，就可以放开种种疑虑，反而可以集中思想把书写得更清楚易读一些。

然而对读者们也有一个要求，即要求他们曾读过一本比较认真的传统的微积分教本。例如同济大学应用数学教研室编的《高等数学》；此书立论平正、平易近人、易教易学，作为进一步学习的出发点是够用的。我对现在的教材绝大部分是肯定的，因为它们帮助读者了解一门科学的基本内容。没有这一点，“重温”二字从何谈起？温故而知新，是希望读者能由此再进一步。这是本书写作的目的。因为假设读者有了这样的基础，我就可以略过许多材料不讲，可以假设读者自己会证明——至少知道——许多定理，特别是本书可以不按顺序，全由读者的兴趣从任何一章开始。这本书不是教本，完全不必要全部读懂。有许多内容则是因自己学力所限没有涉及，特别是有关计算机和数值计算的问题。但是关于本书材料的选择，还有几点深感不尽妥当。一是本书没有一章讲常微分方程，而这是应该有的。这是因为2001年冬我曾应邀在天津大学和南开大学的刘徽应用数学中心讲了一个常微分方程课，学时16，读者和教学目的均与本书相近。但是因为当时还未打算写这本书，那本讲义的内容与本书有些互相牵扯。再重写则头脑里有了一个比较固定的框架，不太好办了，只好等以后有机会再说了。二是本书涉及不少线性代数知识，而基本的又是对偶空间的知识。目前的线性代数教材从一般的线性空间与线性算子角度来讨论的比较少，本书又未能如同处理微积分的古典内容那样细致地处理它，估计会使读者感到困难。如果能在第六、七两章适当补充，效果会更好些。

可是，正如一些数学家一再强调的，数学是“算”懂的，而不是“看”懂的。毫无疑问，这是当前数学教学一大弱点，而且似有加剧之势。本书对所讲的内容力求给以比较清晰的陈述，并给出证明。力求避免“显而易见”、“不难知道”之类的说法。但是对如何有助于读者的计算能力总是考虑不够。一些结论的证明因为与常见的书不同，难免有错，盼读者随时指正。

本书写作首先要感谢老朋友康宏逵老师，他关于数学基础、逻辑以及对微积分的发展的一般看法，多年来与我常共同切磋。本书中引用拉卡托斯的论述更是得益于他。刘伟安老师对大学生（不止是数学专业）的数学教学与我经常讨论。本书中关于向量与张量的处理得他之益甚多。如何给大学生讲一点广义函数更是我们近来交往的主题。田谷基老师对量子力学有兴趣，本书中关于量子力学知识的介绍实际上是我们共同工作的结果。王维克老师从写作本书意念的萌生，一直给予支持和鼓励。全书的打印，刘、田二位老师花了不少精力，对于他们以及许多没有提到的同志，谨致诚挚的谢意。但是我不能不提到责任编辑郭思旭同志。我们是二十多年的老朋友了，我们能合作到退休以后，想起来确有些激动。同样，我向从事编辑和校对的同志深切致谢。

\enlargethispage{1.5\baselineskip}
\begin{flushright}
\begin{minipage}{0.22\textwidth}
\raggedleft
齐民友\\[1mm]
2003年3月
\end{minipage}
\end{flushright}

\clearpage
\tableofcontents
